% ===================================================================
% ML-05c: Musterlösung - Das Verb gustar
% Abgeleitet aus LE-05c (Abschnitt 9: Lösungen)
% ===================================================================

\documentclass[11pt, a4paper]{article}

\usepackage[utf8]{inputenc}
\usepackage[T1]{fontenc}
\usepackage[ngerman]{babel}
\usepackage{helvet}
\renewcommand{\familydefault}{\sfdefault}

\usepackage[margin=2cm]{geometry}
\usepackage{enumitem}
\usepackage{tabularx}
\usepackage{booktabs}

\usepackage{xcolor}
\usepackage{tcolorbox}

\usepackage{fancyhdr}
\usepackage{lastpage}

% Farben
\definecolor{spanisch}{HTML}{c41e3a}
\definecolor{grau}{HTML}{888888}
\definecolor{hellgrau}{HTML}{f5f5f5}
\definecolor{orange}{HTML}{b35c00}
\definecolor{loesung}{HTML}{c62828}

\newcommand{\fachfarbe}{spanisch}

% Variablen
\newcommand{\thema}{Das Verb gustar (gefallen/mögen)}
\newcommand{\sequenz}{05}
\newcommand{\block}{c}
\newcommand{\fach}{Spanisch}
\newcommand{\klasse}{7}
\newcommand{\maxpunkte}{20}

% Kopf-/Fußzeile
\pagestyle{fancy}
\fancyhf{}
\renewcommand{\headrulewidth}{0.4pt}
\renewcommand{\footrulewidth}{0.4pt}
\lhead{\fach{} -- Klasse \klasse}
\chead{\textcolor{loesung}{\textbf{ML-\sequenz\block{} (MUSTERLÖSUNG)}}}
\rhead{}
\lfoot{}
\rfoot{Seite \thepage\ von \pageref{LastPage}}

% Befehle
\newcommand{\punktebox}[1]{\hfill\fbox{\small \_/#1}}
\newcommand{\loesung}[1]{\textcolor{loesung}{\textbf{#1}}}
\newcommand{\luecke}[1]{\rule{#1}{0.4pt}}
\newcommand{\es}[1]{\textcolor{\fachfarbe}{\textbf{#1}}}

% Merkkasten
\newtcolorbox{merkkasten}[1][]{
    colback=hellgrau,
    colframe=\fachfarbe,
    boxrule=1pt,
    arc=0pt,
    left=8pt, right=8pt, top=8pt, bottom=8pt,
    fonttitle=\bfseries\color{\fachfarbe},
    title=#1
}

% Hinweis
\newtcolorbox{hinweis}{
    colback=white,
    colframe=orange,
    boxrule=1pt,
    arc=0pt,
    left=5pt, right=5pt, top=5pt, bottom=5pt,
    fonttitle=\bfseries,
    title=Hinweis
}

% Aufgabe
\newenvironment{aufgabe}[1]{%
    \vspace{10pt}
    \noindent\textbf{\textcolor{\fachfarbe}{Aufgabe #1}}
    \vspace{5pt}
    \begin{list}{}{\leftmargin=0pt}
    \item[]
}{%
    \end{list}
}

% ===================================================================
\begin{document}

% Kopfbereich
\begin{center}
    {\Large\bfseries\textcolor{\fachfarbe}{Arbeitsblatt: \thema}}\\[0.3em]
    {\large\textcolor{loesung}{\textbf{--- MUSTERLÖSUNG ---}}}
\end{center}

\vspace{5pt}

\noindent
\begin{tabularx}{\textwidth}{|X|X|X|}
\hline
\textbf{Name:} \textcolor{loesung}{Musterschüler/in} &
\textbf{Klasse:} \klasse &
\textbf{Datum:} \textcolor{loesung}{XX.XX.XXXX} \\
\hline
\end{tabularx}

\vspace{15pt}

% ===================================================================
% MERKKASTEN (mit Lösungen)
% ===================================================================

\begin{merkkasten}[Merkkasten: Das Verb \textit{gustar} (gefallen/mögen)]

\textbf{Besonderheit:} Im Spanischen sagt man \textit{``Mir gefällt...''} statt \textit{``Ich mag...''}

\vspace{0.5em}

\textbf{Die indirekten Objektpronomen:}
\begin{center}
\begin{tabular}{|l|l|l|}
\hline
\textbf{Pronomen} & \textbf{Deutsch} & \textbf{Beispiel} \\
\hline
me & mir & Me gusta... \\
\hline
te & dir & \loesung{Te gusta...} \\
\hline
le & ihm/ihr/Ihnen & \loesung{Le gusta...} \\
\hline
nos & uns & \loesung{Nos gusta...} \\
\hline
os & euch & \loesung{Os gusta...} \\
\hline
les & ihnen/Ihnen & \loesung{Les gusta...} \\
\hline
\end{tabular}
\end{center}

\vspace{0.5em}

\textbf{Wann gusta, wann gustan?}
\begin{itemize}
    \item \es{gusta} + Infinitiv (Verb): Me gusta \textbf{jugar} al fútbol.
    \item \es{gusta} + Singular-Nomen: Me gusta \textbf{el fútbol}.
    \item \es{gustan} + Plural-Nomen: Me gustan \textbf{los deportes}.
\end{itemize}

\end{merkkasten}

\vspace{10pt}

% ===================================================================
% AUFGABE 1 (mit Lösungen)
% ===================================================================

\begin{aufgabe}{1 -- gusta oder gustan?}
Setze die richtige Form ein: \textit{gusta} oder \textit{gustan}. \punktebox{5}
\end{aufgabe}

\begin{enumerate}[label=\alph*)]
    \item Me \loesung{gusta} el fútbol. \textit{\textcolor{grau}{(Singular: el fútbol)}}
    \item Me \loesung{gustan} los deportes. \textit{\textcolor{grau}{(Plural: los deportes)}}
    \item Te \loesung{gusta} leer libros. \textit{\textcolor{grau}{(Infinitiv: leer)}}
    \item Nos \loesung{gustan} las películas. \textit{\textcolor{grau}{(Plural: las películas)}}
    \item Le \loesung{gusta} la música. \textit{\textcolor{grau}{(Singular: la música)}}
\end{enumerate}

\textit{\textcolor{grau}{Je 1P pro richtige Form}}

\vspace{15pt}

% ===================================================================
% AUFGABE 2 (mit Lösungen)
% ===================================================================

\begin{aufgabe}{2 -- Sätze vervollständigen}
Ergänze die Lücken mit dem richtigen Pronomen und \textit{gusta/gustan}. \punktebox{4}
\end{aufgabe}

\begin{enumerate}[label=\alph*)]
    \item ¿\loesung{Te} \loesung{gusta} jugar al fútbol? (dir)
    \item A María \loesung{le} \loesung{gusta} escuchar música. (ihr)
    \item \loesung{Nos} \loesung{gustan} los videojuegos. (uns)
    \item ¿\loesung{Os} \loesung{gusta} bailar? (euch)
\end{enumerate}

\textit{\textcolor{grau}{Je 1P pro richtige Kombination (Pronomen + Form)}}

\vspace{15pt}

% ===================================================================
% AUFGABE 3 (mit Lösungen)
% ===================================================================

\begin{aufgabe}{3 -- Dialog führen}
Vervollständige den Dialog. Nutze die Redemittel. \punktebox{5}
\end{aufgabe}

\noindent
\textbf{A:} \es{¡Hola!} \loesung{¿Qué te gusta hacer?}

\textbf{B:} \loesung{Me gusta jugar al fútbol.} \es{¿Y a ti?}

\textbf{A:} \es{A mí me gusta escuchar música.} \loesung{¿Te gustan los videojuegos?}

\textbf{B:} \loesung{Sí, me gustan mucho. / No, no me gustan.}

\textbf{A:} \es{¡Genial!}

\vspace{0.5em}
\textit{\textcolor{grau}{Bewertung: Je 1P pro sinnvolle Ergänzung, max. 5P. Alternative Antworten möglich!}}

\vspace{15pt}

% ===================================================================
% AUFGABE 4 (Beispiellösung)
% ===================================================================

\begin{aufgabe}{4 -- Über dich schreiben}
Schreibe 5 Sätze über deine Vorlieben. Benutze verschiedene Pronomen und Aktivitäten. \punktebox{6}
\end{aufgabe}

\textbf{Beispiellösung:}

\vspace{0.5em}
\noindent 1. \loesung{Me gusta escuchar música.}

\vspace{0.5em}
\noindent 2. \loesung{Me gustan los videojuegos.}

\vspace{0.5em}
\noindent 3. \loesung{A mi amigo le gusta jugar al fútbol.}

\vspace{0.5em}
\noindent 4. \loesung{Nos gusta ver películas.}

\vspace{0.5em}
\noindent 5. \loesung{No me gusta bailar.}

\vspace{0.5em}
\textit{\textcolor{grau}{Bewertung:}}
\begin{itemize}
    \item Korrekte gustar-Konstruktion: 3P
    \item Richtige Wahl gusta/gustan: 2P
    \item Verwendung verschiedener Pronomen/Negation: 1P
\end{itemize}

\vspace{20pt}
\hrule
\vspace{10pt}

\begin{flushright}
\textbf{Punkte:} \loesung{\maxpunkte} / \maxpunkte
\end{flushright}

\end{document}
